\documentclass[10pt,a4paper]{amsart}
\usepackage[margin=19mm]{geometry}
\usepackage{amsmath,amssymb,mathtools}
\usepackage[scr=rsfs,cal=euler]{mathalfa}
\usepackage{xfrac}
\usepackage[unicode,hidelinks,bookmarks=false]{hyperref}
\newcommand{\ie}{i.e.\ }
\newcommand{\cf}{cf.\ }
\newcommand{\resp}{respectively\ }
\newcommand{\Subsection}{\subsection}
\providecommand{\nobreakdash}{\nobreak\hbox{-}}
\setlength{\emergencystretch}{2em}
\multlinegap0pt
\usepackage{amssymb}
\newtheorem{prop}{Proposition}
\title{Topology of Bloch Bands from Cauchy Data}
\author{Didier Felbacq, Emmanuel Rousseau}
\date{Source: arXiv:2606.19141v1 (CC BY 4.0)}
\begin{document}
\maketitle

\noindent\emph{Source and licence.} Topology of Bloch Bands from Cauchy Data. Version arXiv:2606.19141v1 (CC BY 4.0), distributed by arXiv under the Creative Commons Attribution 4.0 International licence, http://creativecommons.org/licenses/by/4.0/, as recorded in the arXiv licence metadata for this exact version and shown on the abstract page https://arxiv.org/abs/2606.19141v1. Full licence notice: \texttt{licenses/arxiv-2606.19141-CC-BY-4.0.txt}.\par\smallskip

\noindent\textit{Editorial scope.} Counted unit: the article's prop prop (source0.tex lines 565--567). The article's complete original proof of it is delivered with it (source0.tex lines 569--588, one \texttt{proof} environment of 20 lines). No mathematics is written, repaired or re-derived: the statement and the proof are the article's own lines, in the article's own notation, and no retained line is substituted or re-flowed.  No further source-local context is needed: every cross-reference of every retained line resolves inside the two delivered spans.  Nothing was omitted from any retained span, so the package carries no omission record.  No retained line contains a citation, so the wrapper carries no bibliography and no bibliography entry.  No retained line contains a figure environment, an image-inclusion command or drawing code, so no image is delivered and the package declares no asset dependency.  Editorial formatting only, in addition: an AMS \texttt{amsart} preamble that restates the article's own declarations and macros used by the retained lines, namely the environments \texttt{prop} on separate counters, together with the standard packages those lines need.  The retained environments print in this document as Proposition 1 in order of appearance, whereas the article prints the same items with the numbering of its own full document.  The complete native source of the article as distributed in the arXiv e-print for this exact version is bundled unchanged as \texttt{sources/203121/source0.tex}.
\begin{prop}
For an inversion-symmetric isolated band, the monodromy satisfies $\rho=s_0s_\pi$.
\end{prop}

\begin{proof}
Choose a lifted Bloch section on the fundamental domain
$[0,\pi]$. At the fixed points of the inversion, the section belongs to the eigenspaces
$$
U(0)\in E_{s_0},
\qquad
U(\pi)\in E_{s_\pi}.
$$
To reconstruct the full Brillouin circle, the interval $[\pi,2\pi]$ is obtained from $[0,\pi]$
through the inversion symmetry. The equivariant continuation therefore applies
the inversion representation once at $q=0$ and once at $q=\pi$. Consequently the lifted section acquires the factor
$s_0s_\pi$ after one complete turn around the circle. Hence
$$
U(q+2\pi)=s_0s_\pi\,U(q),
$$
which proves
$$
\rho=s_0s_\pi .
$$
\end{proof}

\end{document}
